\chapter{Settings}\label{app:settings}

The board keeps its settings in nonvolatile memory. Change a setting with
\cmd{set} \textit{key value} in the console, or in \menu{Plotter settings} in
the desktop app. \cmd{config} shows all values. \cmd{defaults} sets the
factory values.

The column ``When'' tells you when a change is possible
(table~\ref{tab:when}): \textbf{live} (at any time), \textbf{idle} (between
jobs), \textbf{disabled} (with the motors disabled), or \textbf{restart}
(after \cmd{restart}).

\section{Machine}

\begin{longtable}{@{}P{0.2\linewidth}P{0.11\linewidth}P{0.16\linewidth}P{0.1\linewidth}P{0.33\linewidth}@{}}
\caption{Machine settings}\label{tab:set-machine}\\
\toprule
\textbf{Key} & \textbf{Default} & \textbf{Range} & \textbf{When} & \textbf{Function} \\
\midrule
\endfirsthead
\toprule
\textbf{Key} & \textbf{Default} & \textbf{Range} & \textbf{When} & \textbf{Function} \\
\midrule
\endhead
\bottomrule
\endfoot
\key{kinematics} & 0 & 0 or 1 & disabled & 0: Cartesian (\conn{M0} = X, \conn{M1} = Y). 1: CoreXY or H-bot (\conn{M0} = A, \conn{M1} = B). Section~\ref{sec:corexy}. A change sets the origin. \\
\key{steps_rev_x} & 800 & 1--1,000,000 & disabled & \conn{M0} steps per revolution. \\
\key{steps_rev_y} & 800 & 1--1,000,000 & disabled & \conn{M1} steps per revolution. CoreXY: the same as \conn{M0}. \\
\key{mm_rev_x} & 5.0 & 0.001--10,000 & disabled & Cartesian: X travel per motor revolution, mm. \\
\key{mm_rev_y} & 5.0 & 0.001--10,000 & disabled & Cartesian: Y travel per motor revolution, mm. \\
\key{pulley_dia_mm} & 12.7324 & 0.1--1000 & disabled & CoreXY: pulley pitch diameter, mm, both motors. Travel per revolution is $\pi$ $\times$ this value. \\
\key{dir_x} & 1 & 1 or -1 & disabled & Direction of X. -1 reverses it. \\
\key{dir_y} & 1 & 1 or -1 & disabled & Direction of Y. -1 reverses it. \\
\key{swap_xy} & 0 & 0 or 1 & disabled & 1: exchange the X and Y axes. A change sets the origin. \\
\key{width_mm} & 400 & 1--100,000 & idle & Plot area in X (the hard clip), mm. \\
\key{height_mm} & 300 & 1--100,000 & idle & Plot area in Y (the hard clip), mm. \\
\end{longtable}

\section{Motion}

\begin{longtable}{@{}P{0.2\linewidth}P{0.11\linewidth}P{0.16\linewidth}P{0.1\linewidth}P{0.33\linewidth}@{}}
\caption{Motion settings}\label{tab:set-motion}\\
\toprule
\textbf{Key} & \textbf{Default} & \textbf{Range} & \textbf{When} & \textbf{Function} \\
\midrule
\endfirsthead
\toprule
\textbf{Key} & \textbf{Default} & \textbf{Range} & \textbf{When} & \textbf{Function} \\
\midrule
\endhead
\bottomrule
\endfoot
\key{draw_mm_s} & 100 & 0.1--10,000 & live & Pen-down speed, mm/s. \cmd{VS} without a value sets it back to this value. \\
\key{travel_mm_s} & 200 & 0.1--10,000 & live & Pen-up speed, mm/s. \\
\key{max_mm_s} & 250 & 0.1--10,000 & live & Highest speed that \cmd{VS} can set, mm/s. \\
\key{job_speed} & 1 & 0 or 1 & live & 1: \cmd{VS} in a job sets the pen-down speed. 0: always \key{draw_mm_s}. \\
\key{accel_mm_s2} & 2000 & 1--1,000,000 & live & Acceleration along the path, mm/s\textsuperscript{2}. Also limits the speed on curves. \\
\key{corner_mm_s} & 20 & 0--10,000 & live & Largest change of speed of one motor at a corner, mm/s. \\
\key{motor_max_mm_s} & 0 & 0--10,000 & live & Highest speed of a motor, as belt travel in mm/s. 0: no limit. Use it for CoreXY. \\
\end{longtable}

\section{Pen}

\begin{longtable}{@{}P{0.2\linewidth}P{0.11\linewidth}P{0.16\linewidth}P{0.1\linewidth}P{0.33\linewidth}@{}}
\caption{Pen settings}\label{tab:set-pen}\\
\toprule
\textbf{Key} & \textbf{Default} & \textbf{Range} & \textbf{When} & \textbf{Function} \\
\midrule
\endfirsthead
\toprule
\textbf{Key} & \textbf{Default} & \textbf{Range} & \textbf{When} & \textbf{Function} \\
\midrule
\endhead
\bottomrule
\endfoot
\key{pen_mode} & 0 & 0 or 1 & disabled & 0: digital output. 1: servo on \conn{M2}. \\
\key{pen_io} & 0 & 0--5 & disabled & Output for the pen, \conn{IO-0}--\conn{IO-5} (\key{pen_mode}~0). \\
\key{pen_active_high} & 1 & 0 or 1 & disabled & 1: the output is on when the pen is down. \\
\key{pen_down_ms} & 150 & 0--10,000 & live & Settle time after the pen goes down, ms. \\
\key{pen_up_ms} & 100 & 0--10,000 & live & Settle time after the pen goes up, ms. \\
\key{z_up_steps} & 0 & $\pm$100,000,000 & disabled & \conn{M2} pen-up position, steps. \\
\key{z_down_steps} & 800 & $\pm$100,000,000 & disabled & \conn{M2} pen-down position, steps. \\
\key{z_vel_steps} & 20,000 & 1--10,000,000 & disabled & \conn{M2} speed, steps/s. \\
\key{z_accel_steps} & 400,000 & 1--100,000,000 & disabled & \conn{M2} acceleration, steps/s\textsuperscript{2}. \\
\end{longtable}

\section{Safety and Jobs}

\begin{longtable}{@{}P{0.2\linewidth}P{0.11\linewidth}P{0.16\linewidth}P{0.1\linewidth}P{0.33\linewidth}@{}}
\caption{Safety and job settings}\label{tab:set-safety}\\
\toprule
\textbf{Key} & \textbf{Default} & \textbf{Range} & \textbf{When} & \textbf{Function} \\
\midrule
\endfirsthead
\toprule
\textbf{Key} & \textbf{Default} & \textbf{Range} & \textbf{When} & \textbf{Function} \\
\midrule
\endhead
\bottomrule
\endfoot
\key{estop_di6} & 1 & 0--2 & live & 0: off. 1: fault when \conn{DI-6} is low. 2: fault when it is high. \\
\key{auto_enable} & 1 & 0 or 1 & live & 1: enable the motors when a job arrives. \\
\key{hlfb_check} & 1 & 0 or 1 & live & 1: the servos must supply HLFB when the motors enable. \\
\key{optimize} & 3 & 0--7 & live & Sources that the optimizer orders: 1 USB, 2 COM, 4 network. Add the values. \\
\key{ident} & \code{7475A} & text & live & Reply to \cmd{OI}. \code{758x} and \code{759x} (DraftMaster) put plotter 0,0 at the center of the bed. \\
\end{longtable}

\section{Serial Port and Network}

\begin{longtable}{@{}P{0.2\linewidth}P{0.11\linewidth}P{0.16\linewidth}P{0.1\linewidth}P{0.33\linewidth}@{}}
\caption{Serial port and network settings}\label{tab:set-links}\\
\toprule
\textbf{Key} & \textbf{Default} & \textbf{Range} & \textbf{When} & \textbf{Function} \\
\midrule
\endfirsthead
\toprule
\textbf{Key} & \textbf{Default} & \textbf{Range} & \textbf{When} & \textbf{Function} \\
\midrule
\endhead
\bottomrule
\endfoot
\key{com_port} & 0 & 0--2 & restart & HP-GL serial port. 0: off. 1: \conn{COM-0}. 2: \conn{COM-1}. \\
\key{com_baud} & 9600 & 300--1,000,000 & restart & Baud rate of the serial port. \\
\key{com_handshake} & 1 & 0--2 & restart & 0: none. 1: XON/XOFF. 2: RTS/CTS. \\
\key{com_rs232} & 1 & 0 or 1 & restart & 1: RS-232 levels. 0: TTL levels. \\
\key{net_mode} & 0 & 0 or 1 & restart & 0: DHCP. 1: static address. \\
\key{ip} & -- & address & restart & Static IP address. \\
\key{netmask} & -- & address & restart & Static netmask. \\
\key{gateway} & -- & address & restart & Static gateway. \\
\end{longtable}

\begin{note}
\notelabel With DHCP, if the board gets no address, it uses 192.168.0.110.
\end{note}
